\section{Marco Teórico}
\label{ch:marco_teorico}

En esta sección se presentan los conceptos teóricos y antecedentes relacionados con el tema de investigación. Se abordarán los fundamentos teóricos que sustentan el estudio, así como las investigaciones previas que han contribuido al desarrollo del conocimiento en esta área.

\subsection{Inteligencia Artificial y sus aplicaciones en sistemas conversacionales}
\label{sec:inteligencia_artificial}

\subsubsection{Conceptos generales de IA y sistemas conversacionales}
\label{subsec:ia_conceptos}

Se define a la Inteligencia Artificial (IA) como el campo de estudio enfocado en el desarrollo de tecnología capaz de realizar tareas que requieren inteligencia humana. Los sistemas basados en IA están diseñados para imitar la capacidad de aprendizaje y razonamiento que poseen los seres humanos \cite{rivera2023inteligencias}.

La IA se utiliza en diversas áreas del conocimiento, como la robótica, los vehículos autónomos y el reconocimiento facial y de objetos. No obstante, en esta sección se abordarán sus aplicaciones en sistemas conversacionales, tales como asistentes virtuales y chatbots.

Una IA generativa de texto es un tipo de Inteligencia Artificial que puede generar texto a partir de datos existentes, y se encuentra en rápido desarrollo con un amplio potencial de aplicaciones \cite{rivera2023inteligencias}. Ejemplos de ello son GPT, desarrollada por OpenAI, y LLaMa, desarrollada por Meta, ambas ampliamente utilizadas en sistemas conversacionales.

En estos sistemas, uno de los principales desafíos es garantizar la coherencia y veracidad de las respuestas, lo que ha dado lugar al estudio de fenómenos como las alucinaciones en modelos de lenguaje, tema que se abordará en la Sección \ref{sec:alucinaciones_modelos_lenguaje}.

\subsubsection{Percepción de los usuarios frente a los chatbots}
\label{subsec:percepcion_chatbots}

Por otro lado, Fondevila-Gascón et al. \cite{fondevila2024chatbot} realizaron un estudio empírico en España para conocer la percepción de los clientes sobre los chatbots como agentes de atención al cliente. Los resultados muestran que la percepción del usuario depende del contexto. Si la pregunta o solicitud es sencilla, la percepción es positiva; si es compleja, es negativa. También se observa que los factores principales que afectan la percepción son la calidad y el contexto.

Además, se concluye que existe una relación entre la calidad del chatbot y la satisfacción del cliente, y se identifica que el factor más importante para el usuario al interactuar con un chatbot es la efectividad. Los autores también señalan que los usuarios prefieren interactuar con agentes humanos en situaciones complejas, como quejas o reclamaciones.

\subsubsection{Interactividad, humanidad y confianza en chatbots}
\label{subsec:confianza_chatbots}

Por su parte, Ding y Najaf \cite{ding2024interactivity} investigaron el impacto de la interactividad y la percepción de humanidad en la confianza hacia los chatbots en el comercio electrónico. Los autores encontraron que tanto la interactividad como la humanidad percibida tienen un efecto positivo en la confianza del usuario.

También identificaron que la confianza actúa como un mediador entre estos factores y la intención de adoptar chatbots. Además, el estudio mostró que el disfrute percibido modera la relación entre la confianza y la intención de adopción, es decir, cuando el usuario disfruta la interacción, la confianza tiene un mayor impacto en su decisión de usar el chatbot. Estos hallazgos resaltan la importancia de diseñar chatbots que no solo sean interactivos y humanos, sino también agradables para el usuario.

\subsection{Procesamiento de Lenguaje Natural (NLP): NLU y NLG}
\label{sec:procesamiento_lenguaje_natural}

\subsection{Alucinaciones en modelos de lenguaje y estrategias de mitigación}
\label{sec:alucinaciones_modelos_lenguaje}

\subsection{Arquitectura RAG (Retrieval-Augmented Generation) y su implementación en sistemas de recomendación}
\label{sec:arquitectura_rag}

\subsection{Inyección de contexto y system prompting para restricción de respuestas}
\label{sec:inyeccion_contexto}

\subsection{Arquitectura multiempresa (multitenancy) en sistemas de software}
\label{sec:arquitectura_multitenant}

\subsection{Estrategias de caché para optimización de consultas}
\label{sec:estrategias_cache}